\chapter{Social and Institutional Projection}
\label{ch:social_institutional}

\epigraph{Local tractability and global dysfunction are consequences of the same
underlying topology: locally consistent systems can be globally incoherent when
gluing conditions across overlapping domains fail.}{}

\section{Platform Architectures as Constrained Trajectory Systems}

\claimP Digital platform architectures can be analyzed as systems that define
admissible trajectories for user behavior while simultaneously optimizing the
platform's own objective functions. Feed algorithms, engagement metrics, and
recommendation systems are constraint geometries that progressively narrow the
space of reachable behavioral configurations.

The narrowing follows a consistent direction: toward configurations that maximize
the platform's measurable objective while restricting access to configurations
that fail to serve those objectives. A user whose work the algorithm does not
amplify finds their reach diminishing regardless of quality. The admissible
trajectory space contracts around platform-legible behavior. This is the Ellul
constraint (Chapter~\ref{ch:ellul}) operating at institutional scale.

\section{Clip Economies and Provenance Erasure}

\claimP The emergence of clip economies---value extracted from fragments disconnected
from their generative contexts---is provenance erasure at the economic scale. A
clip circulates as an autonomous value-bearing object after its inferential
context has been stripped away. Formally this is destructive compression: a
compression that erases structural invariants necessary for understanding the
original object.

The clip is locally coherent---it makes a point---but its inferential continuity
with the discourse it emerged from has been severed. It can be recombined with
other clips in ways its original context would have ruled inadmissible.
Optimization pressure (engagement, virality, platform legibility) acts on the
clips rather than on the underlying claims. Metric drift follows.

\section{Obstruction Cohomology and Institutional Failure}

\claimC Sheaf-theoretic concepts provide a framework for analyzing why certain
institutional configurations cannot be globally stabilized even when locally
consistent. Local solutions resist extension to global solutions when a
non-trivial obstruction class accumulates across overlapping domains.

\begin{definition}{Katz-Admissibility}
An institutional configuration is \emph{Katz-admissible} if every local
resolution of operational inconsistency extends to a globally consistent resolution
across all overlapping domains. An institution failing Katz-admissibility has a
structural obstruction to coherent operation irreducible to individual actor failure.
\end{definition}

\claimC Obstruction failures cannot be resolved by improving local performance or
replacing individual actors. They require restructuring the overlapping domains
so that the obstruction class vanishes, or accepting that a globally consistent
operation is not achievable with the current domain structure.

\section{Labor Markets as Field Systems}

\claimP Platform labor markets can be reinterpreted as field systems in which
workers navigate admissible trajectory spaces defined by rating systems, search
algorithms, and reputation dynamics. The scalar field $\Phi$ measures accumulated
legitimacy; the vector field $\mathbf{v}$ encodes the flows of opportunity; the
entropy field $S$ measures unresolved uncertainty about which trajectories will
remain admissible.

\emph{Participation without guarantee} describes agents whose trajectories are
constrained by systems they cannot inspect or modify. The structural residue of
participation (rating, review, work history) shapes future admissibility under
rules that remain opaque to the participant.

\section{Recursive Productivity Escalation and the Compression Illusion}
\label{sec:ai_labor}

\claimP A recurring pattern in the adoption of productivity technologies is that
labor-saving tools become expectation-amplifying ones. Email did not reduce
communicative obligation; it multiplied it. Spreadsheets did not reduce
analytical labor; they increased demands for continuous forecasting, reporting,
and optimization. Generative AI appears to follow the same pattern at
substantially higher velocity.

The mechanism is a specific instance of proxy stabilization. When AI systems
compress the visible cost of first-pass generation---reducing the time required
to produce a draft, a summary, a code segment, or a design---organizations
frequently reinterpret this compression not as an opportunity to reduce total
labor but as evidence that workers should now produce more outputs in the same
time. If one report previously took a week and now takes a day, the organizational
response is rarely to allocate six fewer working days. It is to expect six reports.

Formally, this can be represented through the projection operator introduced in
the proxy stabilization analysis. Let $X$ be the full operational trajectory
space and $M$ the compressed managerial representation. The operator:
\[
  \pi: X \to M
\]
maps full operational complexity to the compressed representation that
organizational metrics track. AI appears to reduce the cost of operating in $M$,
but the irreducible complexity of $X$ persists. When organizations optimize only
over $M$, invisible reconciliation labor accumulates in the gap $X \setminus
\pi^{-1}(M)$: the substrate labor that maintains coherence between the compressed
map and operational reality.

\claimP This creates what might be called \emph{recursive productivity escalation}.
AI compresses the time required for first-pass generation but simultaneously
increases the amount of oversight, validation, coordination, editing, auditing,
prompting, contextualization, and exception handling required. Workers become
supervisors of probabilistic systems. The labor does not disappear; it mutates.
The bottleneck shifts upward.

\subsection{The Epistemic Inversion}

There is a subtler structural transformation occurring beneath the productivity
narrative. As AI systems produce more text, code, and recommendations,
organizations increasingly require humans to spend cognitive energy verifying
outputs rather than generating them. The worker transitions from creator to
validator. This is more exhausting in a specific sense: error detection under
conditions of high fluency is cognitively taxing in a way that generation is
not. Human cognition is poorly optimized for continuous plausibility filtering
across large volumes of machine-generated content that is locally coherent but
may be globally incoherent.

\claimP This is the inferential versus symbolic continuity distinction from
Chapter~\ref{ch:scaffolded} instantiated at the organizational level. AI systems
produce high symbolic continuity---the outputs read fluently, are locally
plausible, and satisfy surface-level quality filters. The human worker's task
becomes maintaining inferential continuity: detecting the points at which symbolic
fluency has decoupled from derivational constraint. That is precisely the task
that the monograph's reconstruction bound (Chapter~\ref{ch:vocabulary}) shows to
be irreducible under compression.

\subsection{Professional Identity and Occupational Trajectory}

\claimP Many forms of expertise historically involved slow interpretive work:
drafting, synthesis, judgment, and iterative refinement. When AI accelerates
visible production, organizations may begin valuing speed and throughput over
deliberation. This destabilizes occupational identities because workers
increasingly perform orchestration and correction rather than the craftsmanship
through which professional competence was previously legible and measurable.

In trajectory-framework terms: the admissible trajectory space for professional
work has been redrawn by the introduction of AI into the workflow. Trajectories
that previously ran through extended periods of deliberate composition are now
inadmissible---too slow relative to revised organizational tempo. Workers must
navigate a narrowed admissible space whose constraint geometry they did not choose
and cannot easily inspect.

\subsection{Systemic Complexity Growth}

\claimP At the organizational level this produces a structural paradox. AI appears
to increase efficiency locally while increasing systemic complexity globally. Teams
produce more artifacts more rapidly, but coordination overhead rises. More
documentation exists to review. More outputs require alignment. More generated
material must be tracked, versioned, audited, and interpreted. The informational
surface area expands faster than human attention capacity.

This is the Yarncrawler problem at organizational scale: the seam penalty
$L_{\rm seam}$ grows as the number of locally generated artifacts increases,
because each artifact must be glued into a coherent global section with all the
others. AI accelerates local expert production (lowering the cost of each
$\varphi_i$) while the global gluing operation remains bounded by human attention.
The system produces more tiles more quickly but does not increase the rate at
which seam inconsistencies can be repaired.

\claimP Under competitive conditions, productivity gains tend not toward leisure
but toward intensified competition, accelerated expectations, and increased
surveillance granularity. AI becomes a mechanism for tightening admissible
response times across organizations. The title of the relevant empirical
work is precise in this regard: AI does not reduce work because work in
organizations is not a fixed quantity of tasks. Organizational work expands
dynamically to absorb available productive capacity---and AI expands that
capacity while leaving the admissibility conditions for rest and reduction
unchanged.

\begin{constraintboundary}
\textbf{Status:~\textbf{[P]}} The recursive productivity escalation and
epistemic inversion claims are phenomenological interpretations of empirical
organizational research. The formal structure (the $\pi: X \to M$ compression
operator and Yarncrawler seam penalty) provides a modeling framework, not a
derived prediction.

\textbf{Not claimed:} That AI invariably intensifies work in all organizational
contexts, or that the pattern is irreversible. Organizations that deliberately
restructure admissibility conditions around reduced throughput expectations
could break the escalation cycle.

\textbf{Falsification:} Documented cases of AI adoption accompanied by sustained,
organization-wide reduction in total labor hours without corresponding expansion
of output expectations would weaken the escalation claim. Such cases would
constitute constraint-boundary modifications rather than exceptions to the
mechanism.
\end{constraintboundary}

The Spherepop operator Refuse has a social register. Structural
refusal---deliberate non-participation in a system whose admissibility conditions
one rejects---is not merely withdrawal. It commits the refusing agent to a
different trajectory space by eliminating trajectories that run through the
refused system.

\claimP Refusal is most powerful when collective and when it produces structural
residue: when the pattern of refusal becomes legible to other agents and begins
functioning as a signal that alternative admissibility conditions are possible.
\emph{Selection without design} describes the emergence of social order through
locally constrained trajectory filtering without centralized intentional
architecture---the social analogue of the ecological Yarncrawler.

\begin{constraintboundary}
\textbf{Established:} Formal reinterpretations of platform architectures, clip
economies, and labor markets in the trajectory-admissibility framework.
Katz-admissibility as a formal criterion for structural institutional failure.

\textbf{Not established:} That obstruction cohomology is the correct technical
apparatus rather than a productive formal analogy. Institutional applications
carry claim-status \textbf{[P]} or \textbf{[C]} throughout.

\textbf{Speculative~[S]:} Precise quantitative \RSVP field modeling of labor
market dynamics.

\textbf{Falsification:} A platform whose rating-based admissibility filtering
demonstrably tracks underlying competence would weaken the proxy stabilization
claim for that domain. An institution achieving global coherence despite a
non-trivial local obstruction class would falsify Katz-admissibility.
\end{constraintboundary}
